% Informe del proyecto CS4052 - Caso 5: KNN en paralelo con MPI.
% Compilar desde esta carpeta:  latexmk -pdf informe.tex
% Los textos en rojo (\pendiente) marcan lo que falta medir o confirmar.
\documentclass[11pt,a4paper]{article}

\usepackage[utf8]{inputenc}
\usepackage[T1]{fontenc}
\usepackage{lmodern}
\usepackage[spanish,es-tabla,es-nodecimaldot]{babel}
\usepackage[margin=2.4cm]{geometry}
\usepackage{amsmath,amssymb}
\usepackage{booktabs,array,tabularx}
\usepackage{enumitem}
\usepackage{xcolor}
\usepackage{caption}
\usepackage[ruled,vlined,linesnumbered,spanish,onelanguage]{algorithm2e}
\usepackage{tikz}
\usetikzlibrary{babel,arrows.meta,positioning,fit,calc}
\usepackage[hidelinks]{hyperref}

\setlist[itemize]{noitemsep,topsep=3pt}
\captionsetup{font=small}
\setlength{\parskip}{3pt}

% Notación
\newcommand{\ntr}{n_{tr}}
\newcommand{\nte}{n_{te}}
\newcommand{\Xtr}{X_{tr}}
\newcommand{\Xte}{X_{te}}
\newcommand{\ytr}{y_{tr}}
\newcommand{\yte}{y_{te}}
\newcommand{\Tinf}{T_\infty}
\newcommand{\tcomm}[1]{T_{comm_{#1}}}
\newcommand{\tcomp}[1]{T_{comp_{#1}}}

% Marca visible de contenido pendiente. Debe desaparecer antes de entregar.
\newcommand{\pendiente}[1]{\textcolor{red!75!black}{\textbf{[PENDIENTE]} #1}}

% Palabras clave del pseudocódigo
\SetKwFor{ParaPar}{for}{\textcolor{red}{\textbf{pardo}}}{}
\SetKwFor{Para}{for}{do}{}
\SetKwIF{Si}{SinoSi}{Sino}{if}{then}{else if}{else}{}
\SetKwInput{Entrada}{Entrada}
\SetKwInput{Salida}{Salida}
\SetKwComment{Paso}{// }{}
\SetCommentSty{texttt}
\SetKw{Retorna}{return}
\SetKwFunction{Bcast}{Broadcast}
\SetKwFunction{Scatter}{Scatter}
\SetKwFunction{Gather}{Gather}
\SetKwFunction{Reduce}{Reduce}

\begin{document}

% ---------------------------------------------------------------- Título
\begin{center}
  {\large \textbf{CS4052 Computación Paralela y Distribuida} -- 2026-II}\\[2pt]
  Proyecto de Laboratorio -- \textbf{Entrega final}\\[6pt]
  {\LARGE Caso 5: KNN en paralelo con MPI (mpi4py)}\\[4pt]
  Clasificación de dígitos (\textit{load\_digits})\\[6pt]
  Profesor: José Fiestas \quad -- \quad Grupo 7
\end{center}

\begin{center}
\begin{tabular}{lc}
  \toprule
  \textbf{Integrante} & \textbf{Participación} \\
  \midrule
  Barzola Estrella, Johar Jared   & 100\,\% \\
  Huaman Casma, Leonardo Fabian   & 100\,\% \\
  Izaguirre Zavaleta, Luis Fernando & 100\,\% \\
  \bottomrule
\end{tabular}\\[3pt]
{\small (100\,\% indica que el integrante contribuyó en forma equitativa al resto)}
\end{center}

% ---------------------------------------------------------------- 1
\section{Introducción}

KNN (\textit{k-nearest neighbors}) es un algoritmo de clasificación supervisada. Para clasificar un punto nuevo, calcula su distancia a todos los puntos de entrenamiento, toma los $k$ más cercanos y le asigna la etiqueta que más se repite entre ellos.

En este proyecto paralelizamos el código secuencial \texttt{knn\_digits\_sec.py} \cite{base}, que clasifica las imágenes del conjunto \textit{load\_digits} de scikit-learn \cite{sklearn}. El conjunto tiene $n=1797$ imágenes de $8\times 8$ píxeles, es decir $d=64$ características, y 10 clases. El código separa los datos con \texttt{train\_test\_split} (20\,\% para prueba, \texttt{random\_state=42}), lo que da $\ntr=1437$ puntos de entrenamiento y $\nte=360$ de prueba, y usa $k=3$. Su precisión (\textit{accuracy}) es $0{,}9833$, que tomamos como referencia.

El trabajo más costoso es calcular las $\ntr\cdot\nte$ distancias, y cada una se puede calcular de forma independiente de las demás. Por eso es un problema naturalmente paralelizable \cite{slides}.

\textbf{Solución planteada.} Usamos el paradigma de memoria distribuida (paso de mensajes), en Python con mpi4py \cite{mpi4py}. El proceso 0 (maestro) tiene los datos al inicio. Siguiendo el DAG del enunciado \cite{enunciado} (Figura~\ref{fig:dag}):
\begin{itemize}
  \item se envían los datos de prueba a todos los procesos (\texttt{comm.bcast});
  \item se reparten los datos de entrenamiento en bloques (\texttt{comm.scatter});
  \item cada proceso calcula las distancias de todos los puntos de prueba a su bloque, las ordena y guarda sus $k$ vecinos más cercanos por punto de prueba;
  \item el maestro recolecta esos candidatos (\texttt{comm.gather}), elige los $k$ más cercanos de todos y aplica la votación por mayoría.
\end{itemize}

Si un punto es uno de los $k$ vecinos más cercanos entre todos los datos de entrenamiento, también lo es dentro de su bloque: de lo contrario, ese bloque tendría al menos $k$ puntos más cercanos, y el punto no podría ser vecino global. Por ello, al reunir los candidatos locales y seleccionar otra vez los $k$ de menor distancia, se recuperan los mismos vecinos del algoritmo secuencial, salvo empates de distancia.

Para preparar este informe usamos un asistente de IA \cite{ia}; en la Bibliografía describimos para qué lo usamos y en qué partes influyó.

\begin{figure}[ht]
\centering
\begin{tikzpicture}[
  nodo/.style={circle,draw,thick,minimum size=14mm,align=center,font=\scriptsize,inner sep=1pt},
  flecha/.style={-{Stealth[length=2.5mm]},thick},
  com/.style={font=\scriptsize,text=blue!70!black},
  cmp/.style={font=\scriptsize,text=red!75!black}]
  \node[nodo] (ini) at (0,3.6) {Datos\\iniciales};
  \node[nodo] (p0) at (-3,0) {dist\\sort\\k-nearest};
  \node[nodo] (p1) at (-1,0) {dist\\sort\\k-nearest};
  \node at (1,0) {$\cdots$};
  \node[nodo] (pp) at (3,0) {dist\\sort\\k-nearest};
  \node[nodo] (fin) at (0,-3.6) {selección\\y votación};
  \foreach \x in {p0,p1,pp} {
    \draw[flecha] (ini) -- (\x);
    \draw[flecha] (\x) -- (fin);
  }
  \node[com,anchor=east] at (-2.0,2.45) {Scatter: $O\big(\log p\,(a+\tfrac{\ntr}{p}\, b)\big)$};
  \node[com,anchor=west] at (2.0,2.45) {Broadcast: $O\big(\log p\,(a+\nte\, b)\big)$};
  \node[cmp,anchor=east] at (-3.9,0) {$O(\ntr\,\nte/p)$};
  \node[com,anchor=east] at (-2.0,-2.45) {Gather: $O\big(\log p\,(a+\nte\,k\,b)\big)$};
  \node[cmp,anchor=west] at (1.0,-3.6) {$O(\nte\,p\,k)$};
\end{tikzpicture}
\caption{DAG del KNN en paralelo, adaptado de \cite{enunciado} con los costos de \texttt{scatter} y \texttt{gather} corregidos a $O(\log p)$ (Sección~\ref{sec:metricas}), confirmado por el profesor \cite{slides,correo-profesor}. En azul, comunicación; en rojo, cómputo. Aquí $\nte$ cuenta puntos, con $d$ constante.}
\label{fig:dag}
\end{figure}

% ---------------------------------------------------------------- 2
\section{Método}

\subsection{Algoritmo secuencial}
\label{sec:secuencial}

Llamamos $n=\ntr+\nte$ al tamaño de los datos. La división 80/20 se mantiene fija, así que $\nte=c\,n$ y $\ntr=(1-c)\,n$ con $c\approx 0{,}2$. Con esto, el problema depende de un \textbf{solo tamaño} $n$, y $\ntr\,\nte=c(1-c)\,n^2\approx 0{,}16\,n^2$. El tiempo secuencial se obtiene del código base \cite{base} (Tabla~\ref{tab:secuencial}).

\begin{table}[ht]
\centering\footnotesize
\begin{tabular}{ll}
  \toprule
  \textbf{Línea del código secuencial} & \textbf{Operaciones} \\
  \midrule
  \texttt{euclidean\_distance(a, b)} & $d$ restas, $d$ productos, $d-1$ sumas y 1 raíz: $O(d)$ \\
  \texttt{distances = [... for x in X\_train]} & $\ntr\cdot O(d)$ \\
  \texttt{np.argsort(distances)[:k]} & $O(\ntr\log\ntr)$ \\
  \texttt{k\_labels = [y\_train[i] for i in k\_indices]} & $O(k)$ \\
  \texttt{Counter(k\_labels).most\_common(1)} & $O(k)$ \\
  \texttt{y\_pred = [knn\_predict(x, ...) for x in X\_test]} & lo anterior se repite $\nte$ veces \\
  \texttt{accuracy = np.mean(y\_pred == y\_test)} & $O(\nte)$, una sola vez \\
  \bottomrule
\end{tabular}
\caption{Operaciones de cada línea del código secuencial \cite{base}.}
\label{tab:secuencial}
\end{table}

\[
T_s(n)=O\big(\nte(\ntr\, d+\ntr\log\ntr+k)+\nte\big)=O\big(n^2(d+\log n)\big)
\]
Como $d=64>\log_2 n\approx 10{,}8$, el cálculo de distancias domina al ordenamiento:
\[
T_s(n)=O(d\,n^2)
\]

\subsection{Diseño con el método de Foster}
\label{sec:foster}

La Figura~\ref{fig:foster} resume los cuatro pasos de Foster \cite{slides}.

\textbf{Particionamiento.} Usamos descomposición del dominio sobre la data de mayor tamaño, la de entrenamiento ($\ntr>\nte$). La tarea primitiva es la distancia entre un punto de prueba y uno de entrenamiento, lo que da $\ntr\cdot\nte=517\,320$ tareas de costo $O(d)$ y de tamaño parecido, más de cuatro órdenes de magnitud por encima de $p=16$ \cite{slides}.

\textbf{Comunicación.} Cada tarea necesita sus dos puntos, y su distancia se usa luego para elegir los $k$ vecinos de ese punto de prueba. Es decir, las tareas de una misma fila se combinan en una reducción. Contando solo las salidas, son del orden de $\ntr\,\nte\approx 5{,}2\times 10^5$ canales hacia la selección. Esta comunicación es \textit{overhead}: no existe en el secuencial.

\textbf{Aglomeración.} Agrupamos en una tarea todas las columnas de un bloque de puntos de entrenamiento. Así, cada proceso recibe solo su bloque y los datos de prueba, y envía $k$ candidatos por punto de prueba en lugar de todas sus distancias. Contando un canal de Broadcast, uno de Scatter y uno de Gather por bloque, quedan del orden de $3p$ canales.

\textbf{Mapeo.} Asignamos un bloque por proceso. Como $\ntr=1437$ no es múltiplo de $p$ para $p=2,4,8,16$, los primeros $\ntr \bmod p$ procesos reciben $\lceil\ntr/p\rceil$ puntos y el resto $\lfloor\ntr/p\rfloor$. Los bloques difieren en un punto como máximo, así que la carga queda balanceada y en el análisis usamos $\ntr/p$ \cite{slides}. Con $p=16$ son unos 48 canales.

\begin{figure}[ht]
\centering
\begin{tikzpicture}[
  scale=0.42,
  tarea/.style={circle,fill=black!60,inner sep=0pt,minimum size=1.6mm},
  titulo/.style={font=\footnotesize\bfseries,anchor=south},
  flecha/.style={-{Stealth[length=2mm]},thick},
  every node/.append style={font=\scriptsize}]
  % 1. Particionamiento
  \begin{scope}
    \foreach \i in {0,...,3} \foreach \j in {0,...,5} \node[tarea] at (\j,\i) {};
    \node[titulo] at (2.5,4.1) {1. Particionamiento};
    \node at (2.5,-1.1) {$\ntr$ (columnas)};
    \node[rotate=90] at (-1.1,1.5) {$\nte$ (filas)};
  \end{scope}
  % 2. Comunicación
  \begin{scope}[xshift=9.3cm]
    \foreach \i in {0,...,3} {
      \draw[blue!70!black,thick] (0,\i) -- (5,\i);
      \foreach \j in {0,...,5} \node[tarea] at (\j,\i) {};
    }
    \node[titulo] at (2.5,4.1) {2. Comunicación};
    \node at (2.5,-1.1) {reducción por fila};
  \end{scope}
  % 3. Aglomeración
  \begin{scope}[xshift=18.6cm]
    \foreach \i in {0,...,3} \foreach \j in {0,...,5} \node[tarea] at (\j,\i) {};
    \foreach \x/\c in {0/orange,2/teal,4/violet}
      \draw[\c,very thick,rounded corners=2pt,fill=\c,fill opacity=0.15]
        (\x-0.4,-0.4) rectangle (\x+1.4,3.4);
    \node[titulo] at (2.5,4.1) {3. Aglomeración};
    \node at (2.5,-1.1) {bloques de $\ntr/p$};
  \end{scope}
  % 4. Mapeo
  \begin{scope}[xshift=27.9cm]
    \foreach \x/\c/\t in {0/orange/$P_0$,2/teal/$P_1$,4/violet/$P_{p-1}$} {
      \draw[\c,very thick,rounded corners=2pt,fill=\c,fill opacity=0.15]
        (\x-0.4,1.4) rectangle (\x+1.4,3.4);
      \node[text opacity=1] at (\x+0.5,2.4) {\t};
      \draw[flecha,black!70] (\x+0.5,1.3) -- (2.5,0.5);
    }
    \node[draw,thick,rounded corners=2pt,inner sep=2pt] at (2.5,0) {votación};
    \node[titulo] at (2.5,4.1) {4. Mapeo};
    \node at (2.5,-1.1) {un bloque por proceso};
  \end{scope}
\end{tikzpicture}
\caption{Método de Foster aplicado al KNN. Cada punto es una distancia entre un punto de prueba y uno de entrenamiento.}
\label{fig:foster}
\end{figure}

El enunciado también menciona repartir los datos de prueba. Esa alternativa (\textit{scatter} de prueba y \textit{broadcast} de entrenamiento) es válida y evita el paso 6 en el maestro. Sin embargo, seguimos la distribución del DAG, que reparte el conjunto de mayor tamaño.

\subsection{PRAM}
\label{sec:pram}

El Algoritmo~\ref{alg:knn} presenta el pseudocódigo con la notación PRAM del curso \cite{slides,pd02}. No depende de $p$: los límites de los bucles se expresan solo con el tamaño del problema, y \textbf{pardo} va en el bucle sobre $\ntr$, que es la dimensión que se distribuye. El número de procesos aparece recién al ejecutar ese \textbf{pardo}, cuando cada proceso toma un bloque de valores de $j$ y aplica la línea 6 solo a ellos; por eso el Gather reúne varias listas de candidatos por punto de prueba. Esto permite leer directamente el trabajo $W$ y el \textit{span} $\Tinf$ (Sección~\ref{sec:span}).

\begin{algorithm}[ht]
\caption{KNN distribuido con Broadcast/Scatter/Gather$(\Xtr,\ytr,\Xte,\yte,k)$}
\label{alg:knn}
\DontPrintSemicolon
\Entrada{datos de entrenamiento $\Xtr$ ($\ntr\times d$) con etiquetas $\ytr$; datos de prueba $\Xte$ ($\nte\times d$) con etiquetas $\yte$; $k$ vecinos}
\Salida{etiquetas predichas $\hat y[1,\dots,\nte]$ y precisión en el proceso raíz}
\Bcast{$\Xte,k,root=0$}\Paso*[r]{paso 1}
\Scatter{$(\Xtr,\ytr),root=0$}\Paso*[r]{paso 2}
\Para{$i\gets 1$ \KwTo $\nte$}{
  \ParaPar{$j\gets 1$ \KwTo $\ntr$}{
    $D[j]\gets\sqrt{\sum_{l=1}^{d}\big(\Xte[i,l]-\Xtr[j,l]\big)^2}$\Paso*[r]{paso 3}
  }
  $L[i]\gets$ los $k$ pares $(D[j],\ytr[j])$ de menor distancia\Paso*[r]{paso 4}
}
\Gather{$L,C,root=0$}\Paso*[r]{paso 5}
\Si{es el proceso raíz}{
  \Para{$i\gets 1$ \KwTo $\nte$}{
    $V\gets$ los $k$ pares de menor distancia en $C[i]$\Paso*[r]{paso 6}
    $\hat y[i]\gets$ etiqueta más frecuente en $V$\;
  }
  $accuracy\gets$ (cantidad de $i$ con $\hat y[i]=\yte[i]$) $/\ \nte$\Paso*[r]{paso 7}
  \Retorna{$\hat y,accuracy$}\;
}
\end{algorithm}

\textbf{Descripción de los pasos:}
\begin{itemize}
  \item \textbf{Paso 1:} el maestro envía los $\nte$ datos de prueba y $k$ a todos los procesos.
  \item \textbf{Paso 2:} el maestro reparte los datos de entrenamiento y sus etiquetas en bloques.
  \item \textbf{Paso 3:} se calculan las distancias de cada punto de prueba a los puntos de entrenamiento. Cada proceso calcula las de su bloque.
  \item \textbf{Paso 4:} por cada punto de prueba, cada proceso ordena las distancias de su bloque y guarda sus $k$ vecinos locales (distancia y etiqueta).
  \item \textbf{Paso 5:} el maestro recolecta en un solo \texttt{Gather} los candidatos de todos los puntos de prueba, para no pagar la latencia $\nte$ veces.
  \item \textbf{Paso 6:} por cada punto de prueba, el maestro elige los $k$ pares de menor distancia entre todos los candidatos y aplica la votación. No se debe votar directamente con todos los candidatos: se votaría con más de $k$ vecinos y el resultado podría diferir del secuencial.
  \item \textbf{Paso 7:} el maestro calcula la precisión comparando $\hat y$ con $\yte$.
\end{itemize}

Si lo vemos como un PRAM de memoria compartida, el algoritmo es CREW \cite{slides}: todos leen $\Xte$ al mismo tiempo y cada uno escribe solo en sus posiciones. En memoria distribuida, la lectura concurrente se reemplaza por el Broadcast del paso 1.

\subsubsection*{Optimización del paso 6}

En el Algoritmo~\ref{alg:knn}, el maestro hace solo todo el paso 6 y recibe todos los candidatos. Ambos costos crecen con $p$. Podemos evitarlo porque elegir los $k$ menores es una operación asociativa: si $A\oplus B$ denota los $k$ pares de menor distancia de $A\cup B$, entonces $(A\oplus B)\oplus C=A\oplus(B\oplus C)$. A diferencia de una suma \cite{slides}, $\oplus$ no cuesta tiempo constante, sino $O(k)$ por punto de prueba; como $k$ es constante, la cota no cambia. Así, la selección se puede hacer como una \textbf{reducción en árbol}, en la que cada nivel fusiona listas de $k$ pares y solo reenvía $k$ pares por punto de prueba. El Algoritmo~\ref{alg:knnred} reemplaza las líneas 7 a 11 del Algoritmo~\ref{alg:knn}; las líneas 12 y 13 no cambian. En mpi4py corresponde a \texttt{comm.reduce} con una operación propia, que para objetos de Python usa una reducción en árbol por defecto \cite{mpi4py}.

\begin{algorithm}[ht]
\caption{Pasos 5 y 6 optimizados con una reducción en árbol}
\label{alg:knnred}
\DontPrintSemicolon
\Reduce{$L,C,op=\oplus,root=0$}\Paso*[r]{pasos 5 y 6}
\Si{es el proceso raíz}{
  \Para{$i\gets 1$ \KwTo $\nte$}{
    $\hat y[i]\gets$ etiqueta más frecuente en $C[i]$\;
  }
}
\end{algorithm}

Conservamos ambas versiones: el Algoritmo~\ref{alg:knn} cumple con el \texttt{comm.gather} que pide el enunciado, y el Algoritmo~\ref{alg:knnred} es la mejora que compararemos con él.

\subsection{Trabajo y \textit{span}}
\label{sec:span}

Al no depender de $p$, el pseudocódigo permite leer el trabajo $W$ y el \textit{span} $\Tinf$, que es el tiempo con procesadores ilimitados, es decir, con un valor de $j$ por procesador. El bucle sobre $i$ es secuencial y el bucle sobre $j$ es paralelo. Por cada $i$, la línea 5 cuesta $O(d)$ en el camino crítico.

\begin{table}[ht]
\centering\small
\begin{tabular}{lll}
  \toprule
  \textbf{Medida} & \textbf{Con \texttt{Gather} (Alg.~\ref{alg:knn})} & \textbf{Con reducción (Alg.~\ref{alg:knnred})} \\
  \midrule
  Trabajo $W(n)$ & $O(d\,n^2)$ & $O(d\,n^2)$ \\
  Paso 3 en el camino crítico & $O(\nte\, d)=O(d\,n)$ & $O(d\,n)$ \\
  Selección en el camino crítico & $O(\nte\,\ntr)=O(n^2)$ & $O(\nte\, k\log\ntr)=O(k\,n\log n)$ \\
  \textit{Span} $\Tinf(n)$ & $O(n^2)$ & $O\big(n(d+k\log n)\big)$ \\
  Paralelismo $W/\Tinf$ & $O(d)$ & $O\big(d\,n/(d+k\log n)\big)$ \\
  \bottomrule
\end{tabular}
\caption{Trabajo y \textit{span} de ambas versiones. Sin reducción, el maestro revisa todos los candidatos, por lo que su trabajo secuencial acota el paralelismo.}
\label{tab:span}
\end{table}

La Tabla~\ref{tab:span} justifica la optimización: con \texttt{Gather}, el paso 6 del maestro queda en el camino crítico y limita el paralelismo a $O(d)$. Con la reducción, el paralelismo crece aproximadamente como $O(n)$.

Con el teorema de Brent, $T_p=O(W/p+\Tinf)$, pasamos del \textit{span} a un tiempo con $p$ procesos. Como Brent no incluye la comunicación, se suma aparte \cite{slides}:
\[
T_p=O\Big(\frac{d\,n^2}{p}+n^2\Big)+T_{comm},
\qquad
T_p'=O\Big(\frac{d\,n^2}{p}+n(d+k\log n)\Big)+T_{comm}
\]
La primera cota es mayor que la de la Sección~\ref{sec:metricas}, y no es una contradicción. Brent acota la simulación del PRAM, donde el maestro revisa $\ntr$ pares por punto de prueba. Al aglomerar en $p$ bloques, cada bloque envía solo $k$ candidatos, y el maestro revisa $p\,k$ pares. Por eso, en la Tabla~\ref{tab:costos}, el paso 6 cuesta $O(\nte\,p\,k)$ y no $O(\nte\,\ntr)$.

El bucle sobre $i$ también es independiente \cite{enunciado}. Con \textbf{pardo} en $i$ y en $j$, y con la reducción, el \textit{span} bajaría a $O(d+k\log n)$. Dejamos \textbf{pardo} solo en $j$ porque es la dimensión que se distribuye.

\subsection{Métricas}
\label{sec:metricas}

Para la comunicación usamos el modelo del curso \cite{slides}:
\[
T_{comm}=f(p)\,(\alpha+x\beta)
\]
donde $\alpha$ es la latencia, $\beta$ el tiempo de transferencia por unidad de dato, $x$ la cantidad de datos y $f(p)$ depende de la topología. En el DAG, $a\equiv\alpha$ y $b\equiv\beta$.

Siguiendo la corrección del profesor \cite{slides,correo-profesor}, tomamos $f(p)=\log p$ para las tres colectivas, es decir, \textbf{ninguna paga $p$ veces la latencia ni $p$ veces el término de datos}:
\begin{itemize}
  \item \textbf{Broadcast:} cada proceso recibe el mensaje completo, de tamaño $x=\nte\, d$, con costo $\log p\,(a+\nte\, d\, b)$.
  \item \textbf{Scatter:} cada proceso recibe solo su bloque, de tamaño $x=\tfrac{\ntr}{p}\,d$, con costo $\log p\,\big(a+\tfrac{\ntr}{p}\,d\,b\big)$.
  \item \textbf{Gather:} el maestro recibe los $k$ candidatos por punto de prueba de cada proceso, $x=\nte\, k$, con costo $\log p\,(a+\nte\, k\, b)$.
\end{itemize}
Un modelo más detallado describe estas colectivas como un árbol binomial de $\log_2 p$ niveles \cite{thakur}: en el \textit{scatter}, cada nivel reenvía la mitad de lo que tiene, y en el \textit{gather} el maestro termina recibiendo físicamente $(p-1)$ bloques en vez de $p$. Esto puede hacer que la comunicación medida crezca algo más rápido que $\log p$, sobre todo en el \textit{gather}; lo retomamos en la Sección~\ref{sec:escalabilidad} si las mediciones lo muestran.

La Tabla~\ref{tab:costos} muestra el costo de cada paso del Algoritmo~\ref{alg:knn}.

\begin{table}[ht]
\centering\small
\begin{tabular}{clll}
  \toprule
  \textbf{Paso} & \textbf{Líneas} & \textbf{Qué se ejecuta} & \textbf{Costo} \\
  \midrule
  1 & 1 & Broadcast de los datos de prueba & $\tcomm{1}=O\big(\log p\,(a+\nte\, d\, b)\big)$ \\
  2 & 2 & Scatter de los bloques & $\tcomm{2}=O\big(\log p\,(a+\tfrac{\ntr}{p}\, d\, b)\big)$ \\
  3 y 4 & 3--6 & Distancias y ordenamiento & $\tcomp{1}=O\big(\tfrac{\nte\,\ntr}{p}(d+\log\tfrac{\ntr}{p})\big)=O\big(\tfrac{d\,n^2}{p}\big)$ \\
  5 & 7 & Gather de $k$ pares por punto & $\tcomm{3}=O\big(\log p\,(a+\nte\, k\, b)\big)$ \\
  6 y 7 & 8--13 & Selección, votación y precisión & $\tcomp{2}=O(\nte\, p\, k)$ \\
  \bottomrule
\end{tabular}
\caption{Costo de cada paso del Algoritmo~\ref{alg:knn}.}
\label{tab:costos}
\end{table}

\textbf{Tiempo de ejecución.} Sumando los pasos y usando $\nte,\ntr=O(n)$; los tres términos $a\log p$ de la Tabla~\ref{tab:costos} se absorben en uno solo, y el término de datos del \textit{scatter}, $\tfrac{n}{p}\,d\,b\log p$, queda dominado por el del \textit{broadcast}, $n\,d\,b\log p$:
\[
T_p(n,p)=O\Big(\frac{d\,n^2}{p}+k\,n\,p+a\log p+b\,n\log p\,(d+k)\Big)
\]
A diferencia de la versión con el costo del \textit{gather} sin corregir, aquí ningún término de comunicación crece como $p$: el único término lineal en $p$, $k\,n\,p$, viene del cómputo del maestro en el paso 6 (revisar $p\,k$ candidatos por punto de prueba), no de la comunicación.

\textbf{Speedup.} Con $S=T_s/T_p$, donde $T_s$ es la misma rutina de distancias vectorizada, en un proceso y sin MPI. El curso pide especificar la versión secuencial considerada \cite{slides}. No usamos el código base \cite{base} porque el \textit{speedup} mezclaría la vectorización con el paralelismo; su tiempo se informa aparte.
\[
S(n,p)=O\left(\frac{d\,n^2}{\dfrac{d\,n^2}{p}+k\,n\,p+a\log p+b\,n\log p\,(d+k)}\right)
\]

\textbf{Eficiencia.} Con $E=S/p=T_s/(p\,T_p)$, y dividiendo entre $d\,n^2$:
\[
E(n,p)=O\left(\frac{1}{1+\dfrac{k\,p^2}{d\,n}+\dfrac{b\,(d+k)\,p\log p}{d\,n}+\dfrac{a\,p\log p}{d\,n^2}}\right)
=O\left(\frac{1}{1+\dfrac{p^2}{n}}\right)
\]
En el último paso tomamos $a$, $b$, $d$ y $k$ como constantes y conservamos el término dominante, porque $p\log p\le p^2$. Este resultado no cambia respecto a cuando el \textit{gather} todavía llevaba $p$ en su término de datos: el término dominante, $k\,p^2/(d\,n)$, viene del cómputo del paso 6, no de la comunicación.

\textbf{Con el paso 6 optimizado.} La reducción reemplaza $\tcomm{3}$ por $O\big(\log p\,(a+\nte\, k\, b)\big)$ y $\tcomp{2}$ por $O(\nte\, k\log p)$, con lo que ya no queda ningún término que crezca linealmente con $p$:
\[
T_p'(n,p)=O\Big(\frac{d\,n^2}{p}+k\,n\log p+a\log p+b\,n\,(d+k)\log p\Big),
\qquad
E'(n,p)=O\left(\frac{1}{1+\dfrac{p\log p}{n}}\right)
\]

\textbf{Isoeficiencia.} Usamos la función de \textit{overhead} $T_0=p\,T_p-T_s$: para mantener $E$ constante, $T_s$ debe crecer como $T_0$, y cada término de $T_0$ da una condición, de la que manda la más restrictiva \cite{slides}. Con $a$, $b$, $d$ y $k$ constantes, $T_0=O(n\,p^2+n\,p\log p+p\log p)$. Comparando cada término con $T_s=O(n^2)$ se obtienen $n\propto p^2$, $n\propto p\log p$ y $n\propto\sqrt{p\log p}$, así que $n$ debe crecer al menos como $p^2$: $n=\Omega(p^2)$. Con la reducción, $T_0'=O(n\,p\log p+p\log p)$ y basta $n=\Omega(p\log p)$, de modo que esa versión es más escalable. Ambas condiciones se validan con las mediciones en la Sección~\ref{sec:escalabilidad}.

\textbf{Interpretación.} Las expresiones $O(\cdot)$ omiten las constantes de tiempo de cada operación. El curso toma $a$ y $b$ del orden de $10^{-7}$--$10^{-9}$\,s como referencia \cite{slides}; los valores de nuestra máquina se miden en la Sección~\ref{sec:entorno}. Para $p$ pequeño esperamos que el tiempo esté dominado por las distancias y que el algoritmo sea \textit{compute bound}. Al aumentar $p$, ese trabajo disminuye como $1/p$, pero los costos del maestro y de la comunicación no, y la eficiencia tiende a bajar.

\textbf{FLOPs.} Por la fórmula de la distancia euclidiana, cada distancia hace $d$ restas, $d$ productos, $d-1$ sumas y 1 raíz, es decir $3d$ operaciones. La raíz cuenta como una operación aunque sea más lenta, porque los FLOP contabilizan igual operaciones distintas \cite{slides}. En la región paralelizable:
\[
F(n)=3\,d\,\ntr\,\nte,\qquad \text{FLOPs}(p)=\frac{F(n)}{T_p}
\]
También mostramos $F(n)/\max_q t^{(q)}_{comp}$, que excluye la comunicación.
Para $n=1797$ se tiene $F=3\cdot 64\cdot 517\,320\approx 9{,}93\times 10^{7}$ operaciones.

\textbf{Normalización.} Para comparar la teoría con las mediciones, reemplazamos cada término de $T_p$ por un coeficiente en segundos que absorbe $d$, $k$, $c$, $a$, $b$ y el costo real de cada operación:
\[
\hat T(n,p)=c_1\frac{n^2}{p}+c_2\,n\,p+c_3\,n\log_2 p+c_4\log_2 p
\]
Los coeficientes $c_i\ge 0$ se ajustan por mínimos cuadrados con las mediciones. Para la versión con reducción, $\hat T'$ no lleva el término $c_2$. Cada coeficiente tiene un significado físico: $c_1$ refleja el tiempo por operación de distancia, $c_2$ el costo de los candidatos en el maestro, $c_3$ el tiempo por dato y $c_4$ la latencia. Por eso comparamos $c_3$ y $c_4$ con los $b$ y $a$ medidos con un ping-pong \cite{slides}. Además ajustamos por separado $T_{comp}$ y $T_{comm}$, que se miden por separado.

\textbf{Número óptimo de procesos.} Derivando respecto de $p$ e igualando a cero (en $\hat T$ despreciamos además los términos logarítmicos, menores que $c_2\,n\,p$):
\[
p^*\approx\sqrt{\frac{c_1\,n}{c_2}}
\qquad\text{y, con reducción,}\qquad
p^{*\prime}\approx\frac{\ln 2\;c_1\,n^2}{c_3\,n+c_4}
\]
En ambos casos son mínimos, porque la segunda derivada es positiva en ese punto: $\hat T''(p^*)=2\,c_1\,n^2/p^{*3}>0$, y $\big(\hat T'\big)''(p^{*\prime})=(c_3\,n+c_4)/(\ln 2\;p^{*\prime 2})>0$. Los compararemos con el punto en que se cruzan las curvas de cómputo y comunicación \cite{slides}.

% ---------------------------------------------------------------- 3
\section{Resultados y análisis}

\pendiente{Esta sección se completa con las mediciones. Ningún valor experimental ha sido medido todavía.}

\subsection{Entorno experimental}
\label{sec:entorno}
\pendiente{Hardware, sistema operativo, versiones de Python, mpi4py, MPI y NumPy; valores de $p$ y $n$; repeticiones por configuración. Valores de $a$ y $b$ medidos con el ping-pong. Explicar qué reloj se usa: $T_p$ y comunicación salen del maestro, y el cómputo es el máximo entre procesos. Para $n<1797$ se usan submuestras estratificadas, y para $n>1797$ se replican las imágenes con ruido entero, por lo que su precisión no es comparable con la del código base. Aclarar que en un solo nodo MPI transmite por memoria compartida, y que se fijó un hilo por proceso en NumPy.}

\subsection{Desarrollo en versiones beta}
\label{sec:betas}
El código se desarrolló en tres versiones (Tabla~\ref{tab:betas}). Cada una se validó contra el código base \cite{base}: para $n=1797$ y $p=1,2,3,4,7,8,16$, las tres dieron exactamente las mismas 360 predicciones, con precisión $0{,}9833$. Los valores $p=3$ y $p=7$ se incluyeron para comprobar bloques de distinto tamaño. La versión secuencial vectorizada que se usa como $T_s$ también dio las mismas predicciones. Los empates de distancia se resuelven por índice global, para que el resultado no dependa del reparto.

\begin{table}[ht]
\centering\small
\begin{tabularx}{\linewidth}{lX}
  \toprule
  \textbf{Versión} & \textbf{Contenido} \\
  \midrule
  Beta 1 & \texttt{comm.bcast}, \texttt{comm.scatter} y \texttt{comm.gather} con bloques desiguales (Algoritmo~\ref{alg:knn}) \\
  Beta 2 & Beta 1 con medición separada de cómputo y comunicación, y conteo de FLOP \\
  Beta 3 & Paso 6 optimizado con \texttt{comm.reduce} y una operación propia (Algoritmo~\ref{alg:knnred}) \\
  \bottomrule
\end{tabularx}
\caption{Versiones beta del código.}
\label{tab:betas}
\end{table}

Durante el desarrollo aparecieron tres problemas que afectaban las mediciones o el costo teórico:
\begin{itemize}
  \item \textbf{Speedup superlineal falso.} La primera medición daba $E>1$. La causa era que \texttt{(X\_b - x) ** 2} crea arreglos temporales en cada iteración \cite{numpy}: en bloques mayores a unos 0,7\,MB, cada distancia costaba unas 4 veces más, lo que penalizaba a los valores pequeños de $p$. Se corrigió reutilizando un mismo arreglo.
  \item \textbf{Espera de arranque.} No usamos barreras para medir, porque agregan \textit{overhead}. Sin ellas, el maestro empezaba a contar mientras otros procesos seguían importando librerías, y esa espera aparecía como comunicación. Se corrigió ejecutando el algoritmo una vez sin medir y midiendo la segunda ejecución.
  \item \textbf{Costo del paso 6.} La selección se implementó ordenando los $p\,k$ candidatos, para desempatar por índice. Esto agrega un factor $\log(p\,k)$ a $\tcomp{2}$ en la Tabla~\ref{tab:costos}, que en este rango de $p$ es pequeño.
\end{itemize}

\subsection{Tiempos de ejecución, cómputo y comunicación}
\pendiente{Gráficas de los tres tiempos frente a $p$ para cada $n$, precisión frente a $n$ con los datos originales y granularidad $G=T_{comp}/T_{comm}$.}

\subsection{Comparación con la complejidad teórica}
\pendiente{Describir cómo se normalizó la expresión teórica (Sección~\ref{sec:metricas}), con los coeficientes $c_i$ obtenidos y sus unidades. Superponer $\hat T$ y $\hat T'$ a las mediciones, mostrar el ajuste separado de cómputo y comunicación, e informar el error. Comparar $c_3$ y $c_4$ con $b$ y $a$.}

\subsection{Número óptimo de procesos}
\pendiente{Valor de $p^*$ estimado con los $c_i$, y comparación con el $p$ de menor tiempo medido y con el cruce de las curvas de cómputo y comunicación.}

\subsection{\textit{Speedup} y FLOPs}
\pendiente{Gráficas de \textit{speedup} y de FLOPs frente a $p$, y tiempo del código base junto a $T_s$.}

\subsection{Escalabilidad}
\label{sec:escalabilidad}
\pendiente{Eficiencia para distintos $n$ y $p$. Validar la isoeficiencia con los datos: medir $E$ en pares $(n,p)$ con $n/p^2$ constante y con $n/(p\log_2 p)$ constante, comprobar si $E$ se mantiene y compararla con la $E$ del modelo normalizado. Usar cada serie como control de la otra versión.}

\subsection{Mejoras a la solución}
\pendiente{Conclusiones de la comparación entre \texttt{Gather} y la reducción, y otras mejoras: \texttt{argpartition} en el paso 4 y colectivas en mayúscula con arreglos de NumPy.}

% ---------------------------------------------------------------- Bibliografía
\renewcommand{\refname}{Bibliografía}
\begin{thebibliography}{9}
\setlength{\itemsep}{3pt}

\bibitem{enunciado} J. Fiestas, \textit{Computación Paralela y Distribuida: Proyecto -- Laboratorio}, 2026-II. Lima, Perú: UTEC, 2026.\\
\textbf{Impacto:} define el caso 5 y el contenido del informe. De aquí tomamos el DAG (Figura~\ref{fig:dag}) y la distribución de datos de la Introducción (\textit{broadcast} de prueba, \textit{scatter} de entrenamiento, \textit{gather} de candidatos).

\bibitem{base} \texttt{knn\_digits\_sec.py}, código secuencial del caso 5, adjunto al enunciado \cite{enunciado}. [Online]. Available: \url{https://drive.google.com/file/d/1wTBSh7bwT1QVvce0pFMv5tcCWsg4-HeK/view} (accessed Sep. 18, 2026).\\
\textbf{Impacto:} es el código que paralelizamos. De él obtuvimos $T_s$ (Sección~\ref{sec:secuencial}), los parámetros ($k=3$, división 80/20, \texttt{random\_state=42}) y la precisión de referencia ($0{,}9833$).

\bibitem{mpi4py} L. Dalcin, \textit{MPI for Python (mpi4py) documentation} [Online]. Available: \url{https://mpi4py.readthedocs.io/en/stable/} (accessed Nov. 4, 2026).\\
\textbf{Impacto:} revisamos cómo funcionan \texttt{comm.bcast}, \texttt{comm.scatter} y \texttt{comm.gather}, y que las reducciones de objetos usan un árbol por defecto (\texttt{mpi4py.rc.fast\_reduce}), lo que sustenta el Algoritmo~\ref{alg:knnred}.

\bibitem{sklearn} F. Pedregosa \textit{et al.}, ``Scikit-learn: Machine Learning in Python,'' \textit{Journal of Machine Learning Research}, vol. 12, pp. 2825--2830, 2011.\\
\textbf{Impacto:} es la librería que usa el código base (\texttt{load\_digits}, \texttt{train\_test\_split}). De ella salen los datos y su división, descritos en la Introducción.

\bibitem{slides} J. Fiestas, diapositivas y apuntes del curso CS4052: U1.1 Fundamentos de paralelismo, U1.2 Fundamentos (Foster, Flynn), U1.3 DAG, U2.1 Diseño en paralelo (\textit{performance}), U2.2 PRAM 01, U2.3 PRAM 02, U2.4 Particionamiento, U3.2, U3.5, U6.1 y U6.2 (\textit{performance}). Lima, Perú: UTEC, 2026.\\
\textbf{Impacto:} son la base teórica del informe: paralelización natural (Introducción), pasos de Foster y particiones balanceadas (Sección~\ref{sec:foster}), formato del PRAM con \textbf{pardo}, pasos numerados y CREW (Sección~\ref{sec:pram}), teorema de Brent (Sección~\ref{sec:span}), y $S=T_s/T_p$, $E=S/p$, isoeficiencia, FLOPs, medición de $a$ y $b$ y el modelo $T_{comm}=f(p)(\alpha+x\beta)$ (Sección~\ref{sec:metricas}).

\bibitem{pd02} J. Fiestas, \textit{CS4052 -- Práctica Dirigida 02 (solución)}. Lima, Perú: UTEC, 2026.\\
\textbf{Impacto:} nos sirvió de modelo para escribir el pseudocódigo con Scatter, Broadcast y Gather (Sección~\ref{sec:pram}) y para expresar la eficiencia en la forma $E=O(1/(1+\dots))$ (Sección~\ref{sec:metricas}).

\bibitem{thakur} R. Thakur, R. Rabenseifner y W. Gropp, ``Optimization of Collective Communication Operations in MPICH,'' \textit{International Journal of High Performance Computing Applications}, vol. 19, no. 1, pp. 49--66, 2005.\\
\textbf{Impacto:} describe el modelo de árbol binomial para las colectivas, que comentamos en la Sección~\ref{sec:metricas} como explicación física de por qué el maestro recibe $(p-1)$ bloques en el \textit{gather}.

\bibitem{correo-profesor} J. Fiestas, respuesta por correo electrónico a consulta del grupo sobre el costo de \texttt{scatter} y \texttt{gather}, 6 de noviembre de 2026 (comunicación personal).\\
\textbf{Impacto:} confirma que el costo correcto es $O(\log p)$ y no $O(p)$, lo que sustenta la Sección~\ref{sec:metricas}, la Tabla~\ref{tab:costos} y la Figura~\ref{fig:dag}.

\bibitem{numpy} C. R. Harris \textit{et al.}, ``Array programming with NumPy,'' \textit{Nature}, vol. 585, pp. 357--362, 2020.\\
\textbf{Impacto:} es la librería con la que se calculan las distancias en el código paralelo. Su manejo de arreglos temporales explicó el problema de rendimiento descrito en la Sección~\ref{sec:betas}.

\bibitem{ia} Anthropic, \textit{Claude Opus 5} y \textit{Claude Opus 5.5} (asistentes de inteligencia artificial, usados con Claude Code), setiembre y noviembre de 2026.\\
\textbf{Uso:} en la entrega parcial, resumió y ordenó el material de clase y redactó el borrador del informe. En esta entrega, redactó en LaTeX la versión corregida según las observaciones del profesor: pseudocódigo y diagramas (Secciones~\ref{sec:foster} y~\ref{sec:pram}), trabajo y \textit{span} (Sección~\ref{sec:span}) y métricas (Sección~\ref{sec:metricas}). También escribió el código de las versiones beta y los scripts de validación, experimentos y gráficas, y detectó los problemas de medición de la Sección~\ref{sec:betas}. Además, comparó el Método con las diapositivas del curso \cite{slides}.\\
\textbf{Impacto:} la mayor parte de la redacción se hizo con esta herramienta, lo que ahorró tiempo. También propuso la reducción en árbol para optimizar el paso 6 (Algoritmo~\ref{alg:knnred}). El grupo es responsable de lo presentado.

\end{thebibliography}

\end{document}
